
\section{What dynamic financing models establish}
For longer-lived debt, lender receipts depend on future issuance, investment,
priority changes and default choices. A candidate policy $\pi$ produces a
price $q^\pi$. If management optimizes while holding an incompatible price
fixed, the resulting policy may not be the one lenders priced. A simultaneous
solution must make the chosen policy optimal under its price and make the
price value that policy's promised and default-state payments. A constrained
contract-price input is a useful alternative for mechanical stress testing,
but its interpretation is conditional on that price and availability.

Benzoni and coauthors make this distinction explicit in a dynamic model with
fixed issuance costs \cite[Section~2, Propositions~1--2, Appendix~A]{litclosebenzoni}.
Their debt quantity is face value, its price is per unit of face, and
shareholders choose default, issuance or inaction. Value matching and smooth
pasting generate a candidate barrier policy; off-policy deviations and a
supermartingale condition are required for equilibrium verification. The
retained model assumes zero default recovery and risk-neutral agents. It
therefore illustrates the policy-price consistency problem, but is not the
positive-recovery solution required by the preceding example.

The alternative managerial model of Lambrecht and Myers places investment,
borrowing, payout and rents in a dated sources-and-uses constraint and imposes
a no-Ponzi restriction \cite[Sections~2--4]{litcloselambrecht}. Its safe-debt
solutions depend strongly on managerial preferences and governance. In its
exponential-utility example, the liquidation covenant protects lenders under
the stated diffusion dynamics. A jump across the boundary, delayed liquidation
or costly asset disposal would require a different recovery analysis. A
closed-form policy derived under those restrictions is a competitor to
consider, not an identified default policy for every issuer.

Dewatripont, Legros and Matthews provide a primary contracting comparison
\cite[Sections~2--6 and Appendix~B]{litclosedewat}. Their environment has observable
but noncontractible effort, limited liability, monotone investor receipts and
efficient renegotiation. Under entrepreneur-offer renegotiation, a debt
contract can dominate a more general contract in the stated pure-strategy
class. Their appendix also supplies cases in which changed bargaining power,
third parties or mixed messages change that conclusion. Even their investor
conversion option need not be an ordinary equity security, as footnote 16
explains. Contract labels cannot replace the exact contingent payments and
exercise rights in the company model.
